\chapter{Provenance}
\label{ch:provenance}
\index{provenance}

A theorem in a running VNB was either proved in that session or
installed from a record of an earlier proof.  Either way it may rest on
supports that were asserted rather than proved.  And the copy of VNB
that holds it was obtained from somewhere, and may have been changed
since.  This chapter describes how each of these is recorded and how it
is checked:
\begin{itemize}
\item the proof certificate, which lets a load install a theorem without
  running its proof (Section~\ref{sec:proof-certificate});
\item the bill, which names the asserted facts a theorem rests on, and
  the stamp of the VNB test (Sections~\ref{sec:bill}
  and~\ref{sec:test-stamp});
\item signed trees, by which a user records that a copy was checked and
  detects any later change to it (Section~\ref{sec:signed-trees});
\item signed releases, by which the author records the files of a
  published version (Section~\ref{sec:signed-releases}).
\end{itemize}
The last two concern the copy of VNB rather than a single theorem: a
certificate is a plain file of the copy, so what a load accepts as
proven is only as sound as the copy it is read from.  This chapter says
what the two signatures are and what they protect.  The commands that
make and check them are in Chapter~\ref{ch:checking}
(Sections~\ref{sec:signed-trees-use} and~\ref{sec:signed-releases-use}).
The word is used in the sense it has for a single theorem, whose
provenance (\texttt{proven}, \texttt{certified}, \texttt{asserted},
\ldots) says where it came from: here the question is asked of the
library and of the copy as well.

\section{The proof certificate}\label{sec:proof-certificate}
\index{certificate!proof}\index{VNB\_CERTIFIED@\texttt{VNB\_CERTIFIED}}
\index{exam}

\begin{description}
\item[Certifies] that the proof of a named theorem ran to \texttt{qed} and
  passed the inference check, on a stated kernel, and what that proof
  depended on.  It is the mechanism by which the
  everyday load installs a theorem without running its proof.
\item[Issued by] a load that ran the proof: the \emph{exam},
  \texttt{VNB\_CERTIFIED=off ./prover --build-band}, which runs every proof
  and rewrites the store; or an ordinary load meeting a theorem with no
  valid record, which proves it and writes one.
\item[Form] one file per proof file, \texttt{certificates/}$\langle$load
  key$\rangle$\texttt{.cert}, one record per \texttt{qed} in installation
  order.  A record holds the theorem's name, the hash of its statement and
  the statement itself; \texttt{cites}, the names the ledger read off the
  script, each with the hash of its installed statement; \texttt{defs},
  the closure of the statement's defined heads through their definitions
  (the bodies of \texttt{def-functoid}s, the axioms of \texttt{def-predicate}s
  and \texttt{def-constant}s; proven theorems are not followed), with their hashes;
  \texttt{uses}, every macete the kernel saw fire in the proof (an accessor
  read by \texttt{slot} among them), and every \texttt{def-functoid} body and
  predicate defining iff a checker read while accepting an inference, with
  the hash of the rewrite rule or definition, the kernel's own list kept
  independent of the ledger's; the oracles the script called and the bill and oracles
  \texttt{qed} announced; and the kernel hash.  A record carries no date and no host.  A hash is
  the MD5 of the canonical text of a formula, with counter-minted names
  relabelled in order of first occurrence so that it does not depend on
  how much was proved before.  The \emph{kernel hash} is the MD5 of the
  source of the files in \texttt{*cert-kernel-files*}
  (\texttt{certificates.scm}): the files that may call
  \texttt{dg-apply-rule!}, the deduction-graph file, \texttt{expressions},
  \texttt{sequents} and the four checker files; tactics, the kit and
  the drivers are not in it, because every inference they produce is
  checked.
\item[Read by] \texttt{sp}, at load, under \texttt{VNB\_CERTIFIED=on}
  (the default) or \texttt{strict}.  A record is valid when the statement
  being started hashes to its statement hash, every cited, defined and
  used name hashes the same in this tree, and the kernel hash is this
  tree's.  On a valid record the theorem is installed with provenance
  \texttt{certified}, its bill recomputed from the recorded citations, and
  the driver's steps are skipped.  Under \texttt{on} a theorem without a
  valid record is proved and its record written; under \texttt{strict}
  nothing is proved, the uncovered theorems are listed and the load exits
  non-zero.  The bill line of such a theorem reads \texttt{certified modulo
  \{\ldots\}} (a record carries no date and no host, so an exam that changes
  nothing rewrites nothing); the load summary counts \texttt{proven} and
  \texttt{certified} separately, and the word \texttt{proven} alone means
  the proof ran in this image.  The page audit covers the proofs that ran.
\item[Invalidated by] a change to the statement, to the statement of a
  cited name, to a definition the statement mentions or a macete the proof
  fired, or to any kernel file, which invalidates every record at once.  A
  record without a \texttt{uses} field is not valid.
\item[Does not] certify what the record cannot see: a functoid reduced by
  \texttt{beta} needs no entry (a \texttt{beta} contracts only a lambda record
  that reached the goal through the statement, a cited theorem or a named
  unfold, each of which is keyed), and the exam is the backstop for what
  the kernel hash does not cover: the tactics and the kit.  A certificate is not a proof:
  it says a proof ran and was checked.  The design is
  \texttt{docs/certificates-2026-09-24.md}.
\end{description}

\section{The bill, and \texttt{proof-rules}}\label{sec:bill}
\index{bill}\index{proof-rules@\texttt{proof-rules}}

Two further things are certified about a completed proof, neither called
a certificate in the code.  The \emph{bill} that \texttt{qed} prints,
\texttt{proven modulo \{\ldots\} [trust: \ldots]}, is the kernel's
certification of the theorem relative to the asserted facts it names
(Section~\ref{sec:proof-debt}): the theorem holds if those do.  The
command \texttt{(proof-rules)} walks the grounded derivation of the
current proof from its root and tallies every kernel inference by its
rule tag; because every tag is one of the fixed set of kernel operations,
this shows that a proof driven by composite tactics rests on nothing but
those operations, independently of the command trace, which is lossy.

\section{The certification stamp of the VNB test}\label{sec:test-stamp}
\index{certificate!VNB test stamp}

The classification pages of Chapter~\ref{ch:classification} report a
result as ``certified by the VNB test on \textit{date}''.  The stamp is
written to the file \texttt{reference/certification.scm} by
\texttt{vnb-test}, one
\texttt{(certify!\ '\textit{name} "\textit{date}")} per theorem the test
re-proved to \texttt{qed}, and read back at the next load.  It certifies a
run of the test, not a proof.  The everyday mechanism is the proof
certificate above: the exam that writes the store is the run that
re-proves the library.

%% =========================================================
\section{Signed trees}
\label{sec:signed-trees}
\index{signed tree}\index{vnb-trust@\texttt{vnb-trust}}\index{vnb-start@\texttt{vnb-start}}

This section says what a signed tree is and what it protects.  The
commands, \texttt{vnb-trust} and \texttt{vnb-start}, the contents of the
signed list, and the procedure for work that must withstand an
adversary are in Section~\ref{sec:signed-trees-use}.

A user obtains VNB as a copy of its repository: a clone from GitHub,
or an archive of it.  That copy is one directory,
\texttt{\textasciitilde/prover}, called the \emph{tree}.  It holds
\begin{itemize}
\item the Scheme source of the kernel, the tactics and the interface,
  and the Emacs files of the workspace;
\item the launchers \texttt{VNB} and \texttt{prover}, which are short
  scripts;
\item the libraries of definitions (\texttt{structure-library/}) and of
  proofs (\texttt{theorem-library/}), and one proof certificate per
  proof file (\texttt{certificates/});
\item the documentation (\texttt{docs/}, \texttt{reference/});
\item the compiled code (\texttt{.com} files) and the \emph{band}, a
  saved image of the loaded library from which VNB starts.  However,
  the repository does not contain them: a user who clones it makes them on
  the first run.  A copy taken from a machine that has built them
  brings them along.
\end{itemize}
There is no stand-alone executable in the tree.  The program that runs
it is MIT Scheme, with GNU Emacs for the workspace; both are installed
on the computer separately and are not part of the tree.  But the
compiled files and the band are machine code that MIT Scheme runs, so
they are as much the program as the source is, and the signature
covers them.

VNB does not rerun every proof each time it starts.  It installs most
theorems from their proof certificates
(Section~\ref{sec:proof-certificate}), which are plain files in the
tree.  So a program, a person, or an accident that changes files in
the tree can change what VNB accepts as proven, and nothing at
start-up would show it.

A \emph{signed tree} makes such a change visible.  Once the user is
satisfied that a copy of VNB is a good one, the user \emph{accepts}
it: the hash of every file of the tree, the documents aside, is
recorded in a list, and the list is signed with the user's own key.  From then on, a copy is
compared with the lists the user has signed before it is used.  The
picture is
\[
  \begin{array}{c}
  \text{you obtain a copy}\\
  \downarrow\\
  \text{you check the copy}\\
  \downarrow\\
  \text{you sign its hashes}\\
  \downarrow\\
  \text{every later use is compared with them.}
  \end{array}
\]

\paragraph{How a user checks that a copy is good.}  Checking is up to
the user.  Signing is not a check: it only records that the user made
one and was satisfied.  Three ways of checking, from the lightest to
the most thorough:
\begin{itemize}
\item checking the signature of a published release
  (Section~\ref{sec:signed-releases}): the copy is then the one its
  author published;
\item reading what changed since the last accepted copy, which the
  comparison lists file by file;
\item rerunning proofs.  A build with \texttt{VNB\_CERTIFIED=off}
  ignores the certificates and runs every proof through the kernel; it
  is called the \emph{exam}, and takes hours for the whole library.  A
  user can also rerun only some files, for instance those that
  changed.
\end{itemize}
Only the third establishes anything about the proofs themselves.  The
first two establish where the files came from.

\paragraph{What it protects, and what it does not.}  The private key
exists only on the user's computer.  Nothing elsewhere can sign a tree
for the user.  So
\begin{itemize}
\item a copy that differs from every tree the user accepted is refused
  at the next \texttt{vnb-start}, with the differing files named;
\item a proof certificate (Section~\ref{sec:proof-certificate}) is not
  signed, and everything in it can be computed from the tree.  Anyone
  who can write the files of a tree can therefore add a statement to a
  library file and write a certificate for it that the certified load
  accepts.  The statement is then installed without its proof, with the
  provenance \texttt{certified} and an empty bill, and a theorem proved
  from it looks proved.  The signature detects the changed files at the
  next \texttt{vnb-start}; the exam, which ignores the certificates and
  runs every proof, fails at the statement, whose proof does not
  prove it;
\item VNB started without \texttt{vnb-start}, for instance from Emacs
  directly or by \texttt{\textasciitilde/prover/VNB}, is not checked,
  and neither is a tree accepted without a check: the certificate above
  passes both;
\item the check is made when VNB starts.  A file that VNB reads later
  in a session, for instance a library file loaded again, is not
  checked again;
\item MIT Scheme and Emacs are outside the tree and are not checked;
  they are as trustworthy as the way they were installed;
\item the web browser is outside the tree too.  VNB opens its home page
  and the reference pages in it.  The pages are documents and the
  signature does not cover them, so a changed page could show a wrong
  statement; the library that VNB loads is checked, the pages that
  describe it are not;
\item while VNB runs, Emacs answers requests from the home page at a
  local address (\texttt{127.0.0.1}, port 8973).  It carries out only a
  fixed list of nine commands, such as starting a proof or opening the
  calculator, and never runs anything it is sent;
\item nothing here defends a computer that an intruder controls: such
  an intruder can change \texttt{vnb-start}, \texttt{vnb-trust} or
  \texttt{ssh-keygen} itself.
\end{itemize}
The signature says ``I accepted this tree.''  It is not a proof, and it
shows no more about the proofs than the check that preceded it.


%% =========================================================
\section{Signed releases}
\label{sec:signed-releases}
\index{signed release}

This section says what a signed release is and what it shows.  How the
author makes one and how anyone checks one are in
Section~\ref{sec:signed-releases-use}.

A signed tree (Section~\ref{sec:signed-trees}) records a check made by
the user.  A \emph{signed release} records one made by the author: it
lets anyone who obtains a copy of VNB check that the files are the
ones the author published.

\paragraph{What a release is.}  A release has a version,
\texttt{v}\textit{YYYY.MM.DD}, and two files committed under
\texttt{releases/}:
\begin{itemize}
\item the \emph{release manifest}: the commit, then the SHA-256 hash of
  every file git tracks at that commit, one line per file, in the order
  \texttt{git ls-files} gives;
\item the author's SSH signature of the manifest, in the namespace
  \texttt{vnb-release}.
\end{itemize}

\paragraph{The author's key.}  The author's public key is published in two places, under two
different accounts, so that taking over one account is not enough to
replace it:
\begin{itemize}
\item on GitHub, in the author's list of signing keys
  (\url{https://api.github.com/users/JavierThayer/ssh_signing_keys});
\item away from GitHub, in the author's Bluesky profile
  (\url{https://bsky.app/profile/javiertf.bsky.social}): its
  fingerprint, as text in the profile description and in the banner
  image.
\end{itemize}
The key is an ed25519 key, added 2026-11-17, with fingerprint
\begin{center}
\texttt{SHA256:pvv9BBY6uAVZrlz0H7G/xub84NareBQos+3Af8jHTvk}
\end{center}
The README and this manual print it too, for convenience.  They are
not further places to check it, since they come from the same
repository as the release.

\paragraph{What it shows, and what it does not.}  A valid signature
protects a user who obtains a copy of the release from
\begin{itemize}
\item a doctored copy: a fake fork, a tampered mirror, a file changed in
  transit;
\item a forged release, since the private key exists only on the
  author's computer;
\item a planted proof certificate, since certificates are files of the
  release.
\end{itemize}
It does not protect from
\begin{itemize}
\item a change after the download, unless the check is made again;
  \texttt{vnb-start}, used with the downloader's own key, makes it at
  every start;
\item a computer that an intruder already controls;
\item a stolen signing key, until the author withdraws it;
\item a takeover of both of the author's accounts, GitHub and Bluesky,
  together.  A takeover of the GitHub account alone can replace the key
  on GitHub, in the README and in this manual at once; the Bluesky
  copy is what exposes it, which is why it must be compared too.
\end{itemize}

Nor does a valid signature show that the code is harmless: VNB is a
program, and running it means trusting its author or reading it.  And
it shows nothing about the proofs.  For those, no trust in the author
is needed: the exam, run by the user, reruns every proof through the
kernel (Section~\ref{sec:signed-trees-use}, ``For work that must
withstand an adversary'').
